\documentclass[11pt]{article}
\usepackage{amsmath,amssymb,amsthm,mathtools,enumitem,booktabs}
\usepackage[margin=1in]{geometry}

\newtheorem{theorem}{Theorem}
\newtheorem{definition}{Definition}
\newtheorem{remark}{Remark}

\newcommand{\Psym}{\mathsf P_\infty}
\newcommand{\Pinf}{P_\infty}
\newcommand{\Uinf}{\mathcal U_\infty}
\newcommand{\Mz}{M_\zeta}
\newcommand{\F}{\mathcal F}
\newcommand{\Ran}{\operatorname{Ran}}
\newcommand{\Blegal}{\mathcal B}

\title{Step 173: Attempt to Derive the Mellin Kernel of \(\mathsf P_\infty\)}
\author{Six Birds Foundations III}
\date{2026-05-15}

\begin{document}
\maketitle

\section{Verdict}

\[
\boxed{\texttt{kernel\_verdict = V\_operational\_identity}.}
\]

No closed elementary two-variable Mellin kernel is obtained.  The attempt does
derive an exact operational identity for the inherited projection by combining
the Step 104 two-cutoff Sonin definition with the classical
Slepian--Pollak PSWF diagonalization of the time-band limiting operator.  The
identity is explicit enough to define the Mellin-side action of
\(\Psym=\Uinf S_\lambda\Uinf^{-1}\) as a convergent projection formula.

\section{Inherited Records}

Step 102 identifies \(P_\infty\) as the archimedean Sonin/prolate projection.

Step 104 gives the raw archimedean model:
\[
H=L^2(\mathbb R)
\]
or the even subspace, with \(P_\lambda\) the spatial cutoff onto
\[
I_\lambda=[-\lambda,\lambda],
\]
and
\[
\widehat P_\lambda=\F^{-1}P_\lambda\F.
\]
The Sonin space and projection are
\[
\mathcal S_\lambda=\ker P_\lambda\cap\ker\widehat P_\lambda,
\qquad
S_\lambda=\operatorname{Proj}_{\mathcal S_\lambda}.
\]
With the unitary Fourier convention
\[
(\F f)(\xi)=(2\pi)^{-1/2}\int_{\mathbb R}f(x)e^{-ix\xi}\,dx,
\]
the two inherited cutoffs have kernels
\[
P_\lambda(x,y)=1_{I_\lambda}(x)\delta(x-y),
\]
and
\[
\widehat P_\lambda(x,y)
=
\frac{1}{2\pi}\int_{-\lambda}^{\lambda}e^{i(x-y)\xi}\,d\xi
=
\frac{\sin \lambda(x-y)}{\pi(x-y)}
\]
with the diagonal value \(\lambda/\pi\).

Step 153 names the Mellin transport:
\[
T_a=\mathcal M_\Gamma J_a\Uinf^{-1},
\]
where \(\Uinf:H_\infty\to\mathcal H_\infty\) is the archimedean
Hardy--Titchmarsh/Mellin realization.  In the raw \(L^2(\mathbb R)\) log
coordinate used by Step 104, the standard Mellin--Plancherel realization is
the Fourier transform
\[
(\Uinf f)(1/2+i\tau)=(2\pi)^{-1/2}\int_{\mathbb R}f(x)e^{-i\tau x}\,dx,
\]
up to the completed factors already separated in Step 153.

Step 154 introduces the Mellin-side commutator notation
\[
C_\ell=(I-\Pinf)M_{m_\ell}\Pinf
\]
and, in completed Mellin variables,
\[
\widehat r_{\ell,w,k}
=(I-\Psym)M_{m_\ell}\Psym T_a^*
\partial_{\overline w}^{\,k}K_a^\Gamma(\cdot,w).
\]

Step 162 records \(P_\infty\) as the archimedean/Sonin projection in the
Burnol/Sonine pulled-evaluator Hilbert context.  Step 170 records
\[
\Psym=\Uinf P_\infty\Uinf^{-1}
\]
and explicitly notes that the inherited records do not provide
\[
(\Psym F)(s)=\int K_\infty(s,s')F(s')\,d\mu(s').
\]
Step 171 reduces the projected zero-jet gate to
\[
L_{\rho,k}(G)=
\langle \Mz G,\Psym y_{\rho,k}\rangle.
\]

\section{T2a. PSWF Spectral Derivation}

Let
\[
P=P_\lambda,\qquad Q=\widehat P_\lambda.
\]
Then
\[
\mathcal S_\lambda=\ker P\cap\ker Q
=(\Ran P+\Ran Q)^\perp,
\]
so
\[
S_\lambda=I-\operatorname{Proj}_{\Ran P+\Ran Q}.
\]
The key point is that \(\Ran P\) is orthogonal to
\[
(I-P)\Ran Q.
\]
Therefore
\[
\Ran P+\Ran Q=\Ran P\oplus (I-P)\Ran Q.
\]
The projection onto \((I-P)\Ran Q\) is
\[
\Pi_{(I-P)QH}
=(I-P)Q\bigl(Q(I-P)Q\bigr)^{-1}Q(I-P)
\]
on \(QH=\Ran Q\).  Hence the exact two-projection identity is
\[
\boxed{
S_\lambda
=
I-P
-
(I-P)Q\bigl(I_{QH}-QPQ\bigr)^{-1}Q(I-P).
}
\tag{1}
\]

Classical Slepian--Pollak theory diagonalizes the time-band limiting operator
\[
QPQ:QH\to QH.
\]
Let \(\phi_n^\lambda\in QH\) be the bandlimited PSWF basis and let
\[
QPQ\,\phi_n^\lambda=\mu_n^\lambda\phi_n^\lambda,
\qquad 0<\mu_n^\lambda<1.
\]
Then
\[
e_n^\lambda
:=
\frac{(I-P)\phi_n^\lambda}{\sqrt{1-\mu_n^\lambda}}
\]
is an orthonormal basis of \((I-P)QH\), and (1) becomes the strong expansion
\[
\boxed{
S_\lambda
=
I-P
-
\sum_{n\ge0}
\frac{|(I-P)\phi_n^\lambda\rangle
      \langle (I-P)\phi_n^\lambda|}
     {1-\mu_n^\lambda}.
}
\tag{2}
\]
This corrects the tempting but wrong ``select eigenvalues near one'' rule.
The Sonin projection is the complement of both cutoffs.  The PSWF eigenvalues
enter through \((1-\mu_n^\lambda)^{-1}\), and no finite eigenvalue selection is
an exact projection.

\section{T2b. Spatial Kernel and Mellin Transport}

Formula (2) gives the spatial distribution kernel
\[
\boxed{
S_\lambda(x,y)
=
\delta(x-y)
-1_{I_\lambda}(x)\delta(x-y)
-
\sum_{n\ge0}
\frac{(1-1_{I_\lambda}(x))\phi_n^\lambda(x)
      \overline{(1-1_{I_\lambda}(y))\phi_n^\lambda(y)}}
     {1-\mu_n^\lambda}.
}
\tag{3}
\]
The series is an operator identity in the strong projection sense.

Transporting by the standard log-Mellin realization gives, for
\(s=1/2+i\tau\) and \(s'=1/2+i\tau'\),
\[
\boxed{
K_\infty^{\mathrm{op}}(s,s')
=
\delta(\tau-\tau')
-
\frac{\sin \lambda(\tau-\tau')}{\pi(\tau-\tau')}
-
\sum_{n\ge0}\Psi_n^\lambda(s)\overline{\Psi_n^\lambda(s')},
}
\tag{4}
\]
where
\[
\Psi_n^\lambda(s)
=
\Uinf\left(\frac{(I-P_\lambda)\phi_n^\lambda}
{\sqrt{1-\mu_n^\lambda}}\right)(s).
\]
Thus
\[
\boxed{
(\Psym F)(s)
=
F(s)
-
\int \frac{\sin \lambda(\tau-\tau')}{\pi(\tau-\tau')}
F(1/2+i\tau')\,d\tau'
-
\sum_{n\ge0}\Psi_n^\lambda(s)
\langle F,\Psi_n^\lambda\rangle .
}
\tag{5}
\]
Equation (5) is the operational Mellin-side identity obtained in this step.
It is not a closed elementary kernel because the general Mellin transforms of
the PSWF tails \((I-P_\lambda)\phi_n^\lambda\) are not simplified here.

\section{T2c. Sonine-Side Check}

Step 153's Sonine-side evaluator formula gives
\[
\Uinf(J_a^*Y^a_{w,k})
=
T_a^*\partial_{\overline w}^{\,k}K_a^\Gamma(\cdot,w),
\]
but \(K_a^\Gamma\) itself is recorded as a projected Sonine kernel:
\[
K_a^\Gamma(\cdot,w)
=
P_{\mathcal L_a^\Gamma}K_a^{\Gamma,\mathrm{amb}}(\cdot,w).
\]
Therefore the Sonine-side route does not produce a simpler independent
kernel.  It agrees with (5): \(\Psym y_{\rho,k}\) is obtained by applying the
PSWF projection identity to the transported evaluator vector.

\section{T2d. Equivalent Operational Identity}

There is also a kernel-free iteration form.  Let
\[
E=I-P_\lambda,\qquad F=I-\widehat P_\lambda.
\]
The von Neumann alternating projection theorem gives
\[
\boxed{
S_\lambda f
=
\lim_{N\to\infty}(EF)^Nf
}
\tag{6}
\]
strongly in \(L^2(\mathbb R)\).  Since \(E\) and \(F\) have explicit kernels,
(6) is a second exact operational identity.  Formula (5) is preferred because
it diagonalizes the noncommuting cross-term through PSWFs.

\section{Verification}

If \(f\in\Ran P_\lambda\), then \(f\perp\mathcal S_\lambda\), so
\[
S_\lambda f=0.
\]
Formula (1) gives the same conclusion because
\[
Q(I-P)f=0,\qquad (I-P)f=0.
\]
If \(f\in\mathcal S_\lambda\), then \(Pf=Qf=0\), so every correction term in
(1) vanishes and \(S_\lambda f=f\).

\section{\texorpdfstring{\(L_{\rho,k}\)}{Lrho k} Data Point}

The operational identity applies to
\[
L_{\rho,k}(G)=\langle \Mz G,\Psym y_{\rho,k}\rangle.
\]
For the legal zero generator
\[
g_0(t)\equiv0,\qquad G_0(s)=\widehat g_0(s)=0,
\]
and for
\[
\rho_1=\frac12+14.134725\ldots i,\qquad k=0,
\]
one obtains the exact data point
\[
\boxed{
L_{\rho_1,0}(G_0)=0.
}
\]
For a nonzero legal generator \(G\), formula (5) gives
\[
L_{\rho,k}(G)
=
\langle \Mz G,y_{\rho,k}\rangle
-\langle \Mz G,\Uinf P_\lambda\Uinf^{-1}y_{\rho,k}\rangle
-\sum_{n\ge0}
\langle \Mz G,\Psi_n^\lambda\rangle
\langle \Psi_n^\lambda,y_{\rho,k}\rangle.
\]
This is now a concrete PSWF-series evaluation problem rather than an
undefined projection symbol.

\section{Nonclaim Boundary}

This step does not prove RH, does not prove \(\Xi_{\mathrm{BC}}=0\), does not
prove full ZI-COV, and does not prove projected jet-surjectivity.  It does not
replace the inherited Sonin/prolate projection by a zeta-zero spectral
projection.  It does not weaken public-shadow non-promotion, finite-window
Calkin blindness, auxiliary-GRH smuggling, incomplete character spectrum
support-only, scalar identity not carrier identity, or ledger-relative Hecke
non-comparability.

\end{document}
